\onecolumn
\begin{tcolorbox}
[title=Prompt for Measuring the Frequency of Different Thinking Operators, colback=white, colframe=black!75!white, breakable]
\begin{spacing}{1.3}
\refstepcounter{tcbtable}
\label{prompt:frequency}

% You are a “Cognitive Operator Analyzer \& Counter”.

% The user will provide a QA pair. Your task consists of TWO PHASES:

% --------------------------------\\
% PHASE 1: DETAILED ANALYSIS\\
% --------------------------------\\
% You should first perform a detailed internal analysis of the Answer (A), identifying and reasoning about the occurrences of different cognitive operators.

% In this phase:\\
% - You may freely analyze, explain, and reason.\\
% - You should identify where and how each operator appears.\\
% - This analysis is allowed to be verbose and explanatory.\\

% --------------------------------\\
% PHASE 2: FINAL OPERATOR COUNT (STRICT OUTPUT)\\
% --------------------------------\\
% After completing the analysis, you MUST output a clean operator count summary in a STRICT format for automated regex parsing.\\

% ========================\\
% OPERATORS TO COUNT\\
% ========================\\

% 1) DECOMP\_PLAN (Decomposition \& Planning Operator)
% Definition:\\
% A single, coherent planning episode in which a high-level goal is decomposed into sub-goals or steps.\\

% Counting rule (IMPORTANT):\\
% - Count ONE (1) DECOMP\_PLAN per overall planning episode.\\
% - Do NOT count individual sub-steps or sub-plans within the same plan separately.\\
% - If multiple independent planning episodes occur at different points in the Answer, count them separately.\\

% 2) CAUSAL\_INFER (Causal Inference Operator)
% Definition:\\
% A coherent causal inference chain that derives conclusions based on prerequisite conditions or results.\\

% Counting rule (IMPORTANT):\\
% - Count ONE (1) CAUSAL\_INFER per causal chain.\\
% - Multiple causal connectors (e.g., “because”, “therefore”, “thus”, “if–then”) that express the SAME causal relationship are counted ONCE.\\
% - Only count again if a new, independent causal chain is introduced.\\

% 3) MONITOR (Monitoring Operator)\\
% Definition:\\
% Evaluation or verification of the current reasoning path, assumptions, or intermediate conclusions.\\

% Counting rule:\\
% - Each monitoring or self-evaluation event counts as one instance.\\

% 4) BACKTRACK (Backtracking / Step-Revision Operator)\\
% Definition:\\
% Any modification, correction, adjustment, or partial rollback of a previously stated reasoning step, plan, or intermediate conclusion.\\

% Counting rule:\\
% - Count ONE BACKTRACK for each step-level revision or correction event.\\
% - A full restart is NOT required; partial or local revisions still count.\\

% 5) REPR\_REFRAME (Representation Reframing Operator)
% Definition:\\
% A change in perspective, formulation, abstraction, or problem representation that enables further progress.\\

% Counting rule:\\
% - Count each distinct reframing or perspective shift as one instance.\\

% ========================\\
% COUNTING RULES\\
% ========================\\

% - Analyze ONLY the Answer (A). Ignore the Question (Q).\\
% - Count each clearly identifiable operator action once.\\
% - Multiple operators may be counted within the same sentence if they represent distinct actions.\\
% - Assign each instance to the most dominant operator type.\\
% - Do NOT include explanations or text in the final count output.\\

% ========================\\
% FINAL OUTPUT FORMAT (MANDATORY)\\
% ========================\\

% Your FINAL output must appear AFTER your analysis and must contain ONLY the following five lines.
% No extra text, no blank lines, no explanations.\\

% COUNT\_DECOMP\_PLAN:<int>\\
% COUNT\_CAUSAL\_INFER:<int>\\
% COUNT\_MONITOR:<int>\\
% COUNT\_BACKTRACK:<int>\\
% COUNT\_REPR\_REFRAME:<int>\\

You are a \textbf{Cognitive Operator Analyzer \& Counter}. The user will provide a QA pair. Your task has \textbf{TWO PHASES}:

\textbf{Phase 1: Detailed Analysis}  
- Analyze the Answer (A) in detail, identifying occurrences of different cognitive operators.  
- You may explain and reason freely.  
- Identify where and how each operator appears.  

\textbf{Phase 2: Final Operator Count (Strict Output)}  
- After analysis, output a clean operator count summary in strict format for automated parsing.  

\textbf{Operators to Count:}

\begin{enumerate}[leftmargin=*, label=\arabic*)]
    \item \textbf{DECOMP\_PLAN (Decomposition \& Planning Operator)}: One coherent planning episode; sub-steps do not count separately. Multiple independent plans count separately.
    \item \textbf{CAUSAL\_INFER (Causal Inference Operator)}: One per causal chain; multiple connectors in the same chain count as one; new independent chains count separately.
    \item \textbf{MONITOR (Monitoring Operator)}: Each evaluation or self-check counts as one instance.
    \item \textbf{BACKTRACK (Backtracking / Step-Revision Operator)}: Each revision or correction of a reasoning step counts once.
    \item \textbf{REPR\_REFRAME (Representation Reframing Operator)}: Each distinct reframing or perspective shift counts once.
\end{enumerate}

\textbf{Counting Rules:}  
- Analyze only the Answer (A), ignore the Question (Q).  
- Count each operator action once.  
- Multiple operators can appear in the same sentence if distinct.  
- Assign each instance to the most dominant operator.  
- Do not include explanations in the final output.  

\textbf{Final Output Format (Mandatory):}  
Your output must contain only the following five lines:

\begin{verbatim}
COUNT_DECOMP_PLAN:<int>
COUNT_CAUSAL_INFER:<int>
COUNT_MONITOR:<int>
COUNT_BACKTRACK:<int>
COUNT_REPR_REFRAME:<int>
\end{verbatim}
\end{spacing}
\end{tcolorbox}